% from main.tex
%%%%%%%%%%  scan idx 131  |  folio 118  %%%%%%%%%%
acycliques en degré \(\leq a-n\); écrivant le morphisme induit sur les suites exactes de cohomologie des triangles précédents, on en déduit que \(TF\to T'F\) induit un isomorphisme \(\Ho^i(TF)\qis\Ho^i(T'F)\) pour \(i\leq a-n\). Comme on peut choisir \(n\) arbitrairement grand négatif, on a gagné.

b) La démonstration est laissée en exercice au lecteur : voir par exemple ([H] II 3.3).

\textbf{Remarque 3.7.1.} On verra plus bas (7.1.2) un autre cas de « compatibilité » de \(\RHom\) avec \(\mathrm Lf^*\).

\textbf{Corollaire 3.7.2.} (cf. EGA \(0_I\) 5.2.7). \emph{Soient \((S,A)\) un topos annelé, \(s\) un point de \(S\) (SGAA VIII 7.8), \(E,F\in\Ob\D^b(_A S)_{\coh}\). Si \(E_s\qis F_s\), il existe un objet \(s\)-ponctué \(U\) de \(S\) tel que \(E|U\qis F|U\).}

\emph{Preuve.} Soient \(\alpha:E_s\to F_s\) et \(\beta:F_s\to E_s\) des isomorphismes réciproques. D'après (2.22 a)), on a
\[
  \Hom(E,F)_s\qis\Hom(E_s,F_s)
  \qquad
  (\text{resp. }\Hom(F,E)_s\qis\Hom(F_s,E_s)).
\]
Il existe donc un «~voisinage~» \(U\) de \(s\) et une section \(u\) (resp. \(v\)) de \(\Hom(E,F)\) (resp. \(\Hom(F,E)\)) au-dessus de \(U\) telle que \(u_s=\alpha\) (resp. \(v_s=\beta\)). Les homomorphismes \(vu:E|U\to E|U\) et \(uv:F|U\to F|U\) induisent l'identité en \(s\), donc sur un \(V\) convenable au-dessus de \(U\), cqfd.

